\documentclass[a4paper,12pt]{article}
\usepackage{graphicx}
\usepackage{caption}
\usepackage{subcaption}
\usepackage{float}
\usepackage[margin=1in]{geometry}
\usepackage{lmodern}
\usepackage{amsmath}
\usepackage{xcolor}
\usepackage{listings}
\usepackage{pgffor}

\lstset{
  breaklines=true,
  postbreak=\mbox{\textcolor{red}{$\hookrightarrow$}\space},
}

\graphicspath{{../../python/perceptron/output/}}

\begin{document}

\title{Relatório perceptron}
\author{Eduardo}
\date{\today}
\maketitle

A figura~\ref{fig:decisao} apresenta a superfície de decisão gerada pelo modelo perceptron. Essa superfície ilustra como o modelo separa as diferentes classes no espaço de entrada, destacando as regiões atribuídas a cada classe com base nos dados de treinamento.

\begin{figure}[H]
  \centering
  \includegraphics[width=0.5\textwidth]{blobs_decision_boundary.png}
  \caption{Superfície de decisão}
  \label{fig:decisao}
\end{figure}

A figura~\ref{fig:errors} mostra a evolução do erro ao longo das épocas durante o treinamento do modelo perceptron.

\begin{figure}[H]
  \centering
  \includegraphics[width=0.5\textwidth]{blobs_errors.png}
  \caption{Erro por época}
  \label{fig:errors}
\end{figure}

\end{document}